\NoteChapter{2}{Chapter 3 弱拓扑}



\NoteSection{1}{最弱拓扑}

在无穷维空间中，闭且有界的集合通常不是紧集.
例如，$\ell^2$ 的标准正交基 $(e_n)$ 两两之间的距离为 $\sqrt{2}$，
因此不含依范数收敛的子列；但对每个连续线性泛函
$f\in(\ell^2)^\star$，都有 $f(e_n)\to0$.
弱拓扑的基本想法正是只用连续线性泛函来观察向量，
以这些数值的收敛代替范数收敛.

\begin{definition}[映射族生成的最弱拓扑]
    设 $X$ 是集合，$I$ 是指标集，
    $(Y_i)_{i\in I}$ 是一族拓扑空间，
    $\phi_i:X\to Y_i$ 是一族映射.
    由 $(\phi_i)_{i\in I}$ 在 $X$ 上生成的\textbf{最弱拓扑}
    （也称初始拓扑），是使每个 $\phi_i$ 都连续的最小拓扑.
    
\end{definition}


这个定义等价于，
它由所有集合$\phi_i^{-1}(V),\quad i\in I,\quad V\subset Y_i\text{为开集}$
所生成.

\begin{proposition}[邻域基]
    在上述最弱拓扑中，点 $x\in X$ 的一个邻域基由集合
    \[
        \bigcap_{j=1}^{m}\phi_{i_j}^{-1}(V_j)
    \]
    构成，其中 $m<\infty$，且 $V_j$ 是
    $Y_{i_j}$ 中点 $\phi_{i_j}(x)$ 的邻域.
\end{proposition}

\begin{proposition}[收敛与连续性的判别]
    设 $X$ 赋予由 $(\phi_i)_{i\in I}$ 生成的最弱拓扑.
    \begin{enumerate}
        \item 序列 $(x_n)$ 收敛于 $x$，当且仅当对每个 $i\in I$，
        $\phi_i(x_n)\to\phi_i(x)$.
        \item 若 $Z$ 是拓扑空间，则映射 $\psi:Z\to X$ 连续，
        当且仅当对每个 $i\in I$，复合映射
        $\phi_i\circ\psi:Z\to Y_i$ 连续.
    \end{enumerate}
\end{proposition}

\begin{proposition}[Hausdorff 性]
    假设每个 $Y_i$ 都是 Hausdorff 空间.
    则 $X$ 上的最弱拓扑是 Hausdorff 的，
    当且仅当映射族 $(\phi_i)_{i\in I}$ 能分离 $X$ 中的点：
    对任意 $x\ne y$，存在 $i\in I$ 使
    $\phi_i(x)\ne\phi_i(y)$.
\end{proposition}

\begin{example}[逐坐标收敛]
    取 $X=\mathbb{R}^{\mathbb{N}}$，并令
    $\phi_k(x)=x_k$，其中 $x=(x_1,x_2,\ldots)$.
    每个 $Y_k=\mathbb{R}$ 都取通常拓扑.
    由所有坐标映射生成的最弱拓扑就是乘积拓扑，
    其中的序列收敛正是逐坐标收敛.

    例如，标准基序列 $e_n$ 逐坐标收敛于 $0$，
    但在通常的 $\ell^2$ 范数下，$\|e_n\|=1$，
    因而不依范数收敛于 $0$.
\end{example}

\NoteSection{2}{弱拓扑}

\begin{definition}[弱拓扑]
    设 $E$ 是 Banach 空间.
    由所有连续线性泛函 $f\in E^\star$ 在 $E$ 上生成的最弱拓扑，
    称为 $E$ 上的\textbf{弱拓扑}，记作
    $\sigma(E,E^\star)$.
\end{definition}

\begin{proposition}[弱邻域基]
    点 $x_0\in E$ 的一个弱邻域基由集合
    \[
        \bigcap_{j=1}^{m}
        \left\{x\in E:
        |f_j(x)-f_j(x_0)|<\varepsilon_j\right\},
        \qquad f_j\in E^\star,\quad \varepsilon_j>0,
    \]
    构成，其中每个邻域只涉及有限多个连续线性泛函.
\end{proposition}

\begin{remarkbox}
    \textbf{Remark}：
    对于弱邻域基内的集合，有一个几何理解：
    每个集合 $\{x:|f_j(x)-f_j(x_0)|<\varepsilon_j\}$
    都是以仿射超平面为中轴的一条“直带”.
    一个基本弱邻域就是有限条直带的交集，
    因此弱拓扑只能同时限制有限多个线性观测方向.
\end{remarkbox}

\begin{definition}[弱收敛]
    设 $x_n,x\in E$.
    若
    \[
        f(x_n)\to f(x)
        \qquad\text{对每个 }f\in E^\star,
    \]
    则称 $x_n$ \textbf{弱收敛}于 $x$，记作
    $x_n\rightharpoonup x$.
\end{definition}

\begin{proposition}[弱拓扑与范数拓扑]
    弱拓扑比范数拓扑粗，并有以下结论：
    \begin{enumerate}
        \item 范数收敛必然推出弱收敛；反之通常不成立.
        \item 弱开集一定是范数开集，弱闭集一定是范数闭集.
        \item 弱收敛序列必为范数有界序列.
        \item 若 $x_n\rightharpoonup x$，则
        \[
            \|x\|\le\liminf_{n\to\infty}\|x_n\|.
        \]
        \item 若 $E$ 是有限维空间，则弱拓扑与范数拓扑相同.
    \end{enumerate}
\end{proposition}

\begin{example}[单位球面]
    在无穷维 Banach 空间中，单位球面
    $S_E=\{x\in E:\|x\|=1\}$ 是范数闭集，
    但一般不是弱闭集；它的弱闭包是单位闭球 $B_E$.
    从邻域的几何形状看，有限条直带的交集不能把原点
    与整个单位球面分离.
\end{example}

\NoteSection{3}{弱拓扑与凸集}

对一般集合而言，范数闭与弱闭并不等价；
凸性使 Hahn--Banach 分离定理能够把两种闭性联系起来.

\begin{theorem}[闭凸集的强闭与弱闭]
    设 $C\subset E$ 是凸集.
    则 $C$ 在范数拓扑下闭，当且仅当
    $C$ 在弱拓扑 $\sigma(E,E^\star)$ 下闭.
\end{theorem}

\begin{remarkbox}
    \textbf{核心思想}：
    若 $x\notin C$ 且 $C$ 是范数闭凸集，
    Hahn--Banach 分离定理给出一个连续线性泛函，
    将 $x$ 与 $C$ 分开.
    由这个泛函可以构造包含 $x$、且与 $C$ 不相交的弱开邻域，
    因而 $C$ 是弱闭的.
\end{remarkbox}

\begin{corollary}[Mazur 推论]
    若 $x_n\rightharpoonup x$，
    则存在由 $x_n$ 的有限凸组合组成的序列 $(y_n)$，使得
    $\|y_n-x\|\to0$.
\end{corollary}

\NoteSection{4}{弱星拓扑与紧性}

\begin{definition}[弱星拓扑]
    设 $E$ 是 Banach 空间.
    对每个 $x\in E$，定义取值映射
    \[
        J_x:E^\star\to\mathbb{R},
        \qquad J_x(f)=f(x).
    \]
    由所有取值映射 $(J_x)_{x\in E}$ 在 $E^\star$ 上生成的最弱拓扑，
    称为 $E^\star$ 上的\textbf{弱星拓扑}，记作
    $\sigma(E^\star,E)$.
\end{definition}

\begin{proposition}[弱星邻域与弱星收敛]
    弱星拓扑具有以下性质：
    \begin{enumerate}
        \item 点 $f_0\in E^\star$ 的一个弱星邻域基由集合
        \[
            \bigcap_{j=1}^{m}
            \left\{f\in E^\star:
            |f(x_j)-f_0(x_j)|<\varepsilon_j\right\},
            \qquad x_j\in E,\quad\varepsilon_j>0,
        \]
        构成.
        \item 序列 $(f_n)$ 弱星收敛于 $f$，当且仅当它逐点收敛：
        \[
            f_n\overset{\star}{\rightharpoonup}f
            \quad\Longleftrightarrow\quad
            f_n(x)\to f(x)
            \quad\text{对每个 }x\in E.
        \]
        \item $E^\star$ 赋予弱星拓扑后是 Hausdorff 空间.
    \end{enumerate}
\end{proposition}

\begin{proposition}[弱拓扑与弱星拓扑的比较]
    $E^\star$ 上的弱拓扑为
    $\sigma(E^\star,E^{\star\star})$，
    弱星拓扑为 $\sigma(E^\star,E)$.
    弱星拓扑一般比弱拓扑更粗，因此
    \[
        f_n\rightharpoonup f
        \quad\Longrightarrow\quad
        f_n\overset{\star}{\rightharpoonup}f.
    \]
    在有限维空间中，范数拓扑、弱拓扑和弱星拓扑彼此相同.
\end{proposition}

\begin{theorem}[Banach--Alaoglu 定理]
    $E^\star$ 的闭单位球
    $B_{E^\star}=\{f\in E^\star:\|f\|_{E^\star}\le1\}$
    在弱星拓扑 $\sigma(E^\star,E)$ 下是紧集.
\end{theorem}

\begin{remarkbox}
    \textbf{核心思想}：
    把泛函 $f$ 看成由所有数值 $f(x)$ 组成的坐标表，
    从而把 $B_{E^\star}$ 嵌入乘积空间 $\mathbb{R}^E$.
    单位球上的弱星拓扑正是相应的乘积拓扑，
    而每个坐标都落在一个有界闭区间内.
    因此可以利用乘积空间的紧性研究 $B_{E^\star}$.
\end{remarkbox}

\NoteSection{5}{自反空间}

\begin{definition}[典范映射与自反空间]
    设 $E$ 是 Banach 空间.
    典范映射 $J:E\to E^{\star\star}$ 定义为
    \[
        (Jx)(f)=f(x),
        \qquad x\in E,\quad f\in E^\star.
    \]
    Hahn--Banach 定理保证
    $\|Jx\|_{E^{\star\star}}=\|x\|_E$，因而 $J$ 是等距嵌入.
    若 $J$ 是满射，则称 $E$ 为\textbf{自反空间}.
\end{definition}

\begin{example}[自反与非自反空间]
    有限维空间、Hilbert 空间以及 $L^p$ 空间
    （$1<p<\infty$）都是自反空间.
    $L^1$、$L^\infty$ 以及一般的 $C(K)$ 空间则不是自反空间.
\end{example}

\begin{theorem}[Kakutani 定理]
    Banach 空间 $E$ 自反，当且仅当闭单位球 $B_E$
    在弱拓扑 $\sigma(E,E^\star)$ 下紧.
\end{theorem}

\NoteSection{6}{可分性与可度量化}

\begin{proposition}[可分性的关系]
    若 $E^\star$ 可分，则 $E$ 可分.
\end{proposition}

\begin{proposition}[单位球上拓扑的可度量化]
    设 $E$ 是 Banach 空间.
    \begin{enumerate}
        \item $E$ 可分，当且仅当 $B_{E^\star}$ 上的相对弱星拓扑
        $\sigma(E^\star,E)$ 可由某个度量产生.
        \item $E^\star$ 可分，当且仅当 $B_E$ 上的相对弱拓扑
        $\sigma(E,E^\star)$ 可由某个度量产生.
    \end{enumerate}
\end{proposition}

\begin{remarkbox}
    \textbf{Remark}：
    “相对拓扑”是把整个空间的开集与所讨论的单位球相交后得到的拓扑.
    若 $E$ 可分，则 Banach--Alaoglu 定理给出的弱星紧单位球还是可度量的，
    因而有界对偶序列可以抽出弱星收敛子列.
\end{remarkbox}

\NoteSection{7}{一致凸性}

\begin{definition}[一致凸空间]
    Banach 空间 $E$ 称为\textbf{一致凸的}，
    若对每个 $\varepsilon>0$，存在 $\delta>0$，
    使得任意 $x,y\in B_E$ 只要满足
    $\|x-y\|_E\ge\varepsilon$，就有
    \[
        \left\|\frac{x+y}{2}\right\|_E\le1-\delta.
    \]
\end{definition}

\begin{remarkbox}
    \textbf{几何理解}：
    单位球内相距一定距离的两点，其中点必须统一地落入单位球内部.
    Hilbert 空间和 $L^p$（$1<p<\infty$）是一致凸的；
    $\ell^1$ 与 $\ell^\infty$ 的单位球含有平边，因此不一致凸.
\end{remarkbox}

\begin{theorem}[Milman--Pettis 定理]
    一致凸 Banach 空间是自反空间.
\end{theorem}

